\documentclass[11pt]{article}

\usepackage{fullpage}
\usepackage{amssymb}
\usepackage{amsmath}
\usepackage{authblk}
\usepackage[parfill]{parskip}
\usepackage[pdftitle={How to Price Credit Index Options},
            pdfauthor={Ian Castleton},
            pdfkeywords={Credit, Option, Swaption}]{hyperref}

\begin{document}

%\setlength{\parindent}{0cm}

\title{Implementing a Credit Index Option Pricer\\
	\large Version 1.1}

\date{10 Sept 2020}
\author{Ian Castleton}
\affil{Quantitative Analytics Group, Mizuho International}
\maketitle

%\tableofcontents

\section{Introduction}

\subsection{Credit Indices}
A Credit Index consists of a portfolio of many credit default swaps. The Markit iTraxx indices typically contain 125 names, any of which may default at a given time. The credit index is periodically reviewed to update the list of names according to specified criteria, for example including only the highest credit worthy names in a particular country or sector. The most liquid credit index contracts typically have a maturity of 5 years. In order to standardize the contract terms and enable a liquid market, a credit index will typically have a standard premium coupon (e.g. 100bps). This coupon will in general be different to the par credit spread. Entering into a credit index contract therefore requires payment of an up-front fee.

\subsection{Single Name Options}
Considering first an option on a single-name credit default swap, it is possible to derive a closed-form pricing formula. For example Hull~\cite{HullCreditOption} modifies the standard Black-Scholes swaption formula by using the risky-annuity measure. A notable feature of the single name option is that if the underlying credit instrument defaults before option maturity, the option will expire worthless. For this reason the single name option is known as a "knock-out" option.

\subsection{Credit Index Options}
A Credit Index Option is an option to enter into a Credit Index contract on a specified future date. The option has European exercise.

Since a credit-index contains many names, the credit index will NOT knock-out, and will continue to exist even if one or more of the names in the index defaults. An option on a credit index is therefore more complex than a single-name option. 

To summarize the key differences:
\begin{enumerate}
	\item An option on a credit index will not knock-out if one of the underlying names defaults.
	
	\item Any unsettled defaults which occur during the life of the option are settled upon option exercise. i.e. the investor can claim any losses which occur in the index if they choose to exercise the option. This can alter the exercise decision of the option.
	
	\item The option strike is adjusted to take into account the fact that the premium coupons are not at par. This is significant because the market convention is to price the credit index using flat spreads.
\end{enumerate}

For these reasons the option cannot easily be priced using a closed-form Black-Scholes swaption formula. Instead, the implementation described here closely follows the whitepaper by Bloomberg~\cite{BloombergCreditIndexOption}. We also refer to parts of the book by Dominic O'Kane~\cite{DominicOKane}. We obtain close agreement with the Bloomberg terminal when configured to use a flat 5 year credit curve.

\section{Pricing Methodology}

In the following we assume unit notional $N = 1$. We value the forward and option at time $t=0$. The option expiry occurs at time $t_E$, and the credit index maturity at some later time $T$.

\subsection{Forward spread}

Because the credit index option is a no-knockout contract, the option holder is entitled to claim compensation for any unsettled losses which occur before option maturity, if the investor chooses to exercise the option. Effectively, the option holder has protection against defaults from the moment they enter into the option contract. However, the premium payments do not begin until after the option is exercised. The forward spread is therefore the par spread for a CDS with a protection leg starting immediately $(t=0)$, and a premium leg starting at option expiry $(t=t_E)$. Thus the forward spread is given by:
\begin{equation}
S(t_E, T ) = \frac{ \mbox{PV}_{\mbox{protection leg}}
	(0, T)} {\mbox{Annuity}(t_E,T)} \,.
\label{indexForwardSpread}
\end{equation}

Given this value of $S(t_E,T)$, the forward price of the credit index contract is given by
\begin{equation}
V = \xi \left( S(t_E,T)  - c\right) L(t_E,T )
\label{CDScontract}
\end{equation}
where $\xi$ is the payer/receiver indicator, $c$ is the coupon rate, and $L$ the forward risky annuity.


\subsection{Calibrating to the forward spread}
The credit index spread is assumed to follow a log-normal process, namely
\begin{equation}
S(T) = m \exp\left( -\textstyle\frac{1}{2}\sigma^2T + \sigma\sqrt T Z\right)
\label{logNormalSpread}
\end{equation}
where $m$ is a calibration factor which determines the mean of the distribution and $Z$ is a sample from the standard normal distribution. The goal of the calibration process is to find the value of $m$ which causes the log-normal distribution to reproduce the forward contract price from equation~\ref{CDScontract}.

Within the calibration and option valuation steps, it is important to note that Bloomberg use risky annuities which are calculated using flat spreads: For a given spread, the hazard rate is given by
\begin{equation}
\lambda = \frac{S}{1-R}
\end{equation}
where $R$ is the recovery rate. The survival factor $Q$ which is applied to each cashflow of the premium leg risky annuity is then calculated using
\begin{equation}
Q(T) = \exp( -\lambda T)\, .
\end{equation}

The significance of the flat spread is subtle and easily missed, as it means the credit index price is no longer a linear function of spread $S$ (and the calibration factor $m$), since the risky annuity is itself a function of $S(t_E,T; m)$. Writing this dependence out explicitly, the credit index forward price at option expiry is given by:
\begin{eqnarray}
V &=& \xi \left( S(t_E,T)  - c\right) L(t_E,T; S(t_E,T)) \\
\mbox{Fwd PV} = \mathbb{E}_0 \left[ V  \right] &=& \mathbb{E}_0\left[ \xi \left( S(t_E,T)  - c\right) L(t_E,T; S(t_E,T)) \right] \,.
\label{CalibrationExpectation}
\end{eqnarray}

The expectation in equation~\ref{CalibrationExpectation} is evaluated by integrating over the standard normal distribution variable $Z$. The credit index value is given by
\begin{equation}
\mbox{Fwd PV} = \mathbb{E}_0 \left[ V \right] = \int \displaylimits_{-\infty}^{+\infty} 
	\xi \left( S(t_E,T)  - c\right) L(t_E,T; S(t_E,T)) \,
	\phi(Z) \, dZ
\label{calibrationIntegral}	
\end{equation}
where $\phi(Z)$ is the standard normal probability density function.

\subsection{Option Valuation}

The strike of the option contract is adjusted to reflect the fact that underlying credit index has premium coupon payments of $c$ which in general are not equal to the par-spread. The adjusted strike formula is also known as the "up-front fee" calculation, and is given by
\begin{equation}
H(K) = \xi ( c - K )  \, A(K) \, / \, Q_{t_E} \,.
\label{adjustedStrike}
\end{equation}
As with the credit index contract price, the risky annuity $A(K)$ is calculated with flat spreads, in this case at a spread of $K$. Note however that this annuity is scaled by a survival probability $Q_{t_E}$. This is not mentioned in the Bloomberg document~\cite{BloombergCreditIndexOption}, however it is included in the discussion by O'Kane~\cite{DominicOKane} on p209. O'Kane divides through by the survival factor to capture the fact that the exercise price is risk-free until option expiry date, since it is based on the original notional. Empirically, including this adjustment enables closer agreement with the Bloomberg terminal.

With the strike adjustment, the option valuation is given by
\begin{equation}
\mathcal{O} = P_{t_E} \mathbb{E}_0 \left[ ( V_{t_E} + H(K) + D )^+ \right] \,.
\end{equation}
where $P_{t_E}$ is the discount factor to option expiry, and $D$ is the value of any realised defaults during the life of the option.

As with the calibration step, the option price expectation is calculated by integrating over the standard normal variable:
\begin{equation}
\mathcal{O} = P_{t_E}\int \displaylimits_{-\infty}^{+\infty} 
\left[ V_{t_E} + H(K) + D \right]^+ \, 
\phi(Z) \, dZ \,.
\label{optionIntegral}
\end{equation}

\subsection{Default Settlement Amount and Loss Factor}

The Bloomberg whitepaper~\cite{BloombergCreditIndexOption} specifies the default settlment amount upon exercise as
\begin{equation}
D = \xi N_0 l
\end{equation}
where $N_0$ is the initial notional of the credit index at the start of the option, and $l$ is the fractional index loss relative to the initial notional. More specifically, if there are $M$ credit instruments in the index (typically $M=125$) and instrument $i$ of the credit index suffers a default during the life of the option, we can write
\begin{equation}
D = \xi N_0 \frac{( 1 - R_i )}{M}
\end{equation}
where $R_i$ is the recovery rate of instrument $i$.

%The quantity $( 1 - R_i )$ is a realised loss from a default that has actually occurred and is known at the option valuation time. It can therefore be an input to the option pricer. 

When written in this form we see that O'Kane~\cite{DominicOKane} gives a similar expression for the default settlement (page 205 and equation 11.2), namely
\begin{equation}
D = \frac{1}{M}\sum_{i=1}^M 1_{\tau_i\le t_E}(1-R_i)\,.
\end{equation}
Here the total default settlement amount for unit notional is given as a sum over all the defaulted instruments. The indicator function $1_{\tau_i\le t_E}$ is equal to 1 if the default time $\tau_i$ of instrument $i$ is less than or equal to the option expiry time $t_E$. O'Kane further splits the defaults into those which occur up to the valuation time $t$, and those which occur strictly after $t$ and before the option expiry time $t_E$ (page 210 and equation 11.7):
\begin{eqnarray}
D &=& \frac{1}{M}\sum_{i=1}^M 1_{0<\tau_i\le t}(1-R_i) \label{defaultsBeforeValuationTime}\\
  &+& \frac{1}{M}\sum_{i=1}^M 1_{t<\tau_i\le t_E}(1-R_i)\label{potentialDefaultsAfterValuationTime} \,.
\end{eqnarray}
The defaults in the first term~\ref{defaultsBeforeValuationTime} are known with certainty at valuation time $t$. The second term captures the value of {\em potential} defaults which are not yet known as of time $t$.

The Bloomberg pricer includes only those defaults which are actually realised during the life of the option contract (i.e. the term in~\ref{defaultsBeforeValuationTime}); it does not include the effect of {\em potential} defaults up to time $t_E$ ( the term in~\ref{potentialDefaultsAfterValuationTime}), other than through the forward spread adjustment (see comments on page 4 of the Bloomberg whitepaper for confirmation).

\section{Implementation}

\subsection{Numerical Calibration Procedure}

A Newton-Raphson solver is used to find the calibration parameter $m$ which reproduces the forward contract price in equation~\ref{calibrationIntegral}. It is desirable to evaluate the calibration integral quickly as it is invoked multiple times by the solver. Since the integrand is of the form
\begin{equation}
f(x) \, \exp(-x^2)
\end{equation}
where $f(x)$ is a smooth function, we choose to use the Gauss-Hermite routine from QuantLib, evaluated over 64 points.

\subsection{Evaluating the Option Expected Payoff}

The option payoff integrand in equation~\ref{optionIntegral} is not smooth because it contains a maximum function $( ... )^+$. Gauss-Hermite is therefore not suitable, and we choose instead to use the Gauss-Legendre integrator from QuantLib. Since the QuantLib implementation integrates over the limits $\pm 1$, we perform a change of variable to modify the integration limits to be suitable for the standard normal variable. Sufficient accuracy is achieved with 400 points.

\subsection{Greeks}

The assumptions used to calculate various greeks were chosen to match the Bloomberg terminal. Since the option value is obtained through a numerical quadrature, the greeks were calculated via a bump-and-reval approach.

\subsubsection{Vega}

For the vega calculation, all market data is held constant and the option implied volatility is bumped up by $b=1\%$. The vega is given as follows:
\begin{equation}
v = \mbox{PV}(\sigma + b) - \mbox{PV}(\sigma) \,.
\end{equation}

The MLIB implementation actually allows the user to specify a different bump size, in which case the vega calculation becomes
\begin{equation}
v = \frac{\mbox{PV}(\sigma + b) - \mbox{PV}(\sigma)}{b} \times 0.01 \,.
\end{equation}
Note that in this case a scaling factor of 0.01 is used to ensure that the calculated vega represents the change in PV for a 1\% move in volatility.

\subsubsection{CS01}

The CS01 greek measures the change in option price with a flat parallel shift in credit spread of $b=1\mbox{bp}$. The credit model $M$ is re-calibrated by applying the flat shift to all CDS calibration spreads ${\bf s}$, and the option PV re-calculated using the bumped credit model:
\begin{equation}
\mbox{CS01} = \mbox{PV}(M({\bf s}+b)) - \mbox{PV}(M({\bf s})) \,.
\end{equation}

\subsubsection{Theta}

Bloomerg define theta as ``Time-decay in option value for a one day decrease in option expiry". To strictly capture the time-decay only and match the Bloomberg theta value, we reduce the time to expiry by one calendar day and hold everything else constant. i.e. the market data, forward spread, annuities and calibration factor are all held constant. The reduced time to expiry $T^\prime$ is used in the payoff integral, and replaces the value of $T$ when calculating the log-normal term in equation~\ref{logNormalSpread}. The theta is then the difference in bumped and unbumped PV:
\begin{equation}
\Theta = \mbox{PV}(T^\prime) - \mbox{PV}(T) \,.
\end{equation}


%\pagebreak
\bibliographystyle{plain}
\bibliography{CreditIndexOptionImplementation}

\end{document}

